\documentclass[../../../main.tex]{subfiles}

\begin{document}

\section{The Gradient}

Finally!! Gradients!!! I am very excited to discuss about the
gradients because it has so much applications outside math.
Physics and machine learning are two famous examples that constantly
uses gradients as their foundation. I mean gradient descent is based
one gradients! Now before we get too excited, let's get started with
the definition.

\begin{definition}{}{}
Given a function $f$ whose first partials at $(x_0, y_0)$ exists,
the \textbf{gradient} of $f$ at the point is defined as the following.
\vspace{-0.8\baselineskip}
\[
\nabla f(x_0, y_0)
\coloneqq \langle f_x(x_0, y_0), f_y(x_0, y_0) \rangle
\]
\end{definition}

When we say this out loud, we say del, nabla, or grad $f$,
where ``grad $f$'' is a popular choice. We can see that we can
write $\nabla f$ as the following as well.
\[
\nabla f
= \langle
    \frac{\partial f}{\partial x}, \frac{\partial f}{\partial y}
\rangle
\]
Now from this definition and our earlier discussion on directional
derivative, I think it is not difficult to see that the following
also holds for a unit vector $\mathbf{u}$.
\[
D_\mathbf{u} f(x_0, y_0) = \nabla f(x_0, y_0) \cdot \mathbf{u}
\]
We can also see that the gradient is actually a function! It takes
in a point and spits out a vector. Also, in that sense $\nabla$ is
merely an operator on functions. With this in mind, here is the
theorem that makes gradients so powerful!

\begin{theorem}{}{}
Let $f(x, y)$ be a function where $\nabla f(x, y) \neq \mathbf{0}$
for all points in the domain. Then, the direction of the maximum
slope is $\nabla f(x, y)$ and the magnitude of the maximum slope is
$\| \nabla f(x, y) \|$.
\end{theorem}
\begin{proof}
First, consider the following equation for a unit vector $\mathbf{u}$.
\[
D_\mathbf{u} f(x_0, y_0)
= \nabla f(x, y) \cdot \mathbf{u}
= \| \nabla f(x, y) \| \| \mathbf{u} \| \cos\theta
\]
Notice that for $D_\mathbf{u} f(x_0, y_0)$ to be greatest,
$\cos\theta = 1$ and $\theta = 0$. In other words,
$D_\mathbf{u} f(x_0, y_0)$ is greatest when $\mathbf{u}$ and
$\nabla f(x, y)$ points towards the same direction and
$D_\mathbf{u} f(x_0, y_0) = \| \nabla f(x, y) \|$.
\end{proof}

Continuing with this theorem and proof, I don't think it is too
difficult to see that the minimum of $D_\mathbf{u} f(x_0, y_0)$ occurs
at the opposite direction of $\nabla f(x, y)$ and its magnitude is
$-\| \nabla f(x, y) \|$.

One side note is when $\nabla f(x, y) = \mathbf{0}$, then this point
is a \textbf{saddle point}. Before we wrap up this section,
here is one important theorem.

\begin{theorem}{}{Level Curve Orthogonal}
Let $f(x, y)$ be a function that is differentiable at $(x_0, y_0)$ and
$\nabla f(x_0, y_0) \neq \mathbf{0}$. Then, $\nabla f(x_0, y_0)$ and
the tangent line that passes through level curve at the point are
orthogonal.
\end{theorem}
\begin{proof}
Let $\mathbf{r} = \langle x(t), y(t) \rangle$ be a smooth
parametrization of the level curve $f(x, y) = c$ for some constant
$c$. Using the chain rule, the following equations hold.
\begin{align*}
    f_x(x(t), y(t)) x'(t) + f_y(x(t), y(t)) y'(t) &= 0 \\
    \nabla f(x(t), y(t)) \cdot \mathbf{r}'(t) &= 0
\end{align*}
Therefore, the theorem holds.
\end{proof}

This is it for gradients! It was a relatively short section although
it is such a powerful tool. The fact that gradients give the minimum
slope of a function is very good because we can use this information
to descend towards the local minimum of loss functions in machine
learning. Although I have so much to talk about here, I will save the
rest for the Machine Learning part of LibreNotebook.

\end{document}


